\section{问题二模型建立和求解}

\subsection{决策变量与约束}

设 FY1 的航向角为 $\theta$，从 $x$ 轴正向逆时针计，水平单位方向向量为
\begin{equation}
    \bm e(\theta)=(\cos\theta,\sin\theta,0)^{\mathsf T}.
\end{equation}
为避免同时把投放点和起爆点作为冗余变量，选取
\begin{equation}
    \bm x=(\theta,v,t_e,\tau)^{\mathsf T}
    \label{eq:q2-decision}
\end{equation}
作为决策变量，其中 $v$ 为无人机速度，$t_e$ 为起爆时刻，$\tau$ 为引信延时。投放时刻为 $t_r=t_e-\tau$，由此反算投放点与起爆点：
\begin{align}
    \bm r_r&=\bm u_0+v\bm e(\theta)(t_e-\tau), \\
    \bm r_e&=\bm u_0+v\bm e(\theta)t_e
    +(0,0,-\tfrac12g\tau^2)^{\mathsf T}.
    \label{eq:q2-points}
\end{align}
设 M1 到达假目标原点的时刻为 $t_M=\lVert\bm m_0\rVert/300$，则可行域为
\begin{equation}
    0\le\theta<2\pi,\quad 70\le v\le140,\quad
    0\le\tau\le t_e\le t_M,\quad
    1800-\tfrac12g\tau^2\ge0.
    \label{eq:q2-constraints}
\end{equation}
起爆后烟幕球心满足
\begin{equation}
    \bm c(t)=\bm r_e-(0,0,3(t-t_e))^{\mathsf T},\qquad
    t_e\le t\le\min\{t_e+20,t_M\}.
\end{equation}

\subsection{优化模型与求解算法}

沿用问题一的完整圆柱最坏视线距离 $D(t)$，将有效时刻集合定义为
\begin{equation}
    \Omega(\bm x)=\left\{t\in[t_e,\min\{t_e+20,t_M\}]:D(t;\bm x)\le10\right\}.
\end{equation}
问题二的优化模型为
\begin{equation}
    \max_{\bm x}\quad T_{\mathrm{eff}}(\bm x)
    =\mu\!\left(\Omega(\bm x)\right),
    \qquad \text{s.t.}\quad \cref{eq:q2-constraints},
    \label{eq:q2-objective}
\end{equation}
其中 $\mu(\cdot)$ 表示时间区间并集的长度。由于大部分试探策略完全未命中视线，直接以遮蔽时长搜索会形成大面积零值平台。程序在尚未形成有效区间时，以烟幕球到最坏视线的最小超距作为连续引导；出现有效区间后，再切换为式\cref{eq:q2-objective} 的真实时长目标。采用三个独立随机种子的差分进化搜索，并对候选解作局部细化；最终以每个圆周 $5760$ 个角向点和 Brent 事件根重新计算。完整流程见\cref{fig:q2-flow}。

\begin{figure}[!htbp]
    \centering
    \includegraphics[width=0.76\textwidth]{fig_flow_q2.pdf}
    \caption{问题二单弹联合优化与可行性复核流程}
    \label{fig:q2-flow}
\end{figure}

\subsection{最优投放策略}

数值搜索得到航向角 $\theta=176.641048^{\circ}$，即相对 $x$ 轴负向偏向 $y$ 轴正向 $3.358952^{\circ}$；飞行速度取允许下界 $70\,\mathrm{m/s}$。最优策略在接令后立即投放，投放时刻 $t_r=0$，引信延时与起爆时刻均为 $2.496794\,\mathrm{s}$。相应速度向量、投放点和起爆点为
\begin{align}
    \bm v_U&=(-69.879744,\ 4.101384,\ 0)^{\mathsf T}\ \mathrm{m/s},\\
    \bm r_r&=(17800.000,\ 0.000,\ 1800.000)^{\mathsf T}\ \mathrm{m},\\
    \bm r_e&=(17625.524700,\ 10.240309,\ 1769.453506)^{\mathsf T}\ \mathrm{m}.
\end{align}
策略汇总见\cref{tab:q2-strategy}。

\begin{table}[!htbp]
    \centering
    \caption{问题二 FY1 单枚烟幕弹优化策略}
    \label{tab:q2-strategy}
    \begin{tabular}{lc}
        \toprule[1.5pt]
        决策量或结果 & 数值 \\
        \midrule[1pt]
        航向角（$+x$ 轴逆时针） & $176.641048^{\circ}$ \\
        FY1 速度 & $70.000000\,\mathrm{m/s}$ \\
        投放时刻与投放点 & $0\,\mathrm{s}$；$(17800.000,0.000,1800.000)\,\mathrm{m}$ \\
        起爆时刻与起爆点 & $2.496794\,\mathrm{s}$；$(17625.525,10.240,1769.454)\,\mathrm{m}$ \\
        完全遮蔽区间 & $[2.496794,7.039674]\,\mathrm{s}$ \\
        最大有效遮蔽时长 & $4.542880\,\mathrm{s}$ \\
        \bottomrule[1.5pt]
    \end{tabular}
\end{table}

完整圆柱从起爆后几乎立即进入完全遮蔽状态，于 $t=7.039674\,\mathrm{s}$ 离开有效烟幕范围，故
\begin{equation}
    T_{\mathrm{eff}}=7.039673606-2.496793631
    =4.542879975\,\mathrm{s}\approx\boxed{4.543\,\mathrm{s}}.
\end{equation}
与问题一的给定策略相比，遮蔽时长增加 $3.151\,\mathrm{s}$，提高约 $226.4\%$。优化收敛过程及最优策略下的距离裕度见\cref{fig:q2-diagnostics}。

\begin{figure}[!htbp]
    \centering
    \includegraphics[width=0.96\textwidth]{q2_optimization_diagnostics.pdf}
    \caption{问题二多初值优化收敛与最优策略距离裕度}
    \label{fig:q2-diagnostics}
\end{figure}

\subsection{收敛性与边界检验}

三个独立随机种子所得时长分别为 $4.542617$、$4.542868$ 和 $4.542834\,\mathrm{s}$，随后在活跃边界上细化至 $4.542880\,\mathrm{s}$。圆周角向采样由 $90$ 加密至 $1440$ 时结果保持 $4.542879993\,\mathrm{s}$，加密至 $5760$ 后仅变化 $1.79\times10^{-8}\,\mathrm{s}$。对最优点作单侧扰动：速度增至 $70.1\,\mathrm{m/s}$ 时，时长降至 $4.542795\,\mathrm{s}$；投放延后 $0.01\,\mathrm{s}$ 时，降至 $4.542390\,\mathrm{s}$。这说明最低速度与立即投放是受题面约束限制的活跃最优边界。所有有效时刻均早于导弹到达假目标，且起爆高度为正，策略满足物理与时序约束。

\endinput
\section{问题二模型建立和求解}

数学建模必然涉及不少数学公式的使用。下面简单介绍一个可能用得上的数学环境。

首先是行内公式，例如 $ \theta $ 是角度。行内公式使用 \verb|$  $| 包裹。

行间公式不需要编号的可以使用 \verb|\[  \]| 包裹，例如
\[
E=mc^2
\]
其中 $ E $ 是能量，$ m $ 是质量，$ c $ 是光速。

如果希望某个公式带编号，并且在后文中引用可以参考下面的写法：
\begin{equation}
E=mc^2
\label{eq:energy}
\end{equation}
式\cref{eq:energy}是质能方程。

多行公式有时候希望能够在特定的位置对齐，以下是其中一种处理方法。
\begin{align}
P & = UI \\
& = I^2R
\end{align}
\verb|&| 是对齐的位置， \verb|&| 可以有多个，但是每行的个数要相同。

矩阵的输入也不难。
\[
\mathbf{X} = \left(
    \begin{array}{cccc}
    x_{11} & x_{12} & \ldots & x_{1n}\\
    x_{21} & x_{22} & \ldots & x_{2n}\\
    \vdots & \vdots & \ddots & \vdots\\
    x_{n1} & x_{n2} & \ldots & x_{nn}\\
    \end{array} \right)
\]

分段函数这些可以用 \verb|case| 环境，但是它要放在数学环境里面。
\[
f(x) =
    \begin{cases}
        0 &  x \text{为无理数} ,\\
        1 &  x \text{为有理数} .
    \end{cases}
\]
在数学环境里面，字体用的是数学字体，一般与正文字体不同。假如要公式里面有个别文字，则需要把这部分放在 \verb|text| 环境里面，即 \verb|\text{文本环境}| 。

公式中个别需要加粗的字母可以用 \verb|$\bm{math symbol}$| 。如 $ \alpha a\bm{\alpha a} $ 。

以上仅简单介绍了基础的使用，对于更复杂的需求，可以阅读相关的宏包手册，如 \href{http://texdoc.net/texmf-dist/doc/latex/amsmath/amsldoc.pdf}{amsmath}。

希腊字母这些如果不熟悉，可以去查找符号文件 \href{http://mirrors.ctan.org/info/symbols/comprehensive/symbols-a4.pdf}{symbols-a4.pdf} ，也可以去 \href{http://detexify.kirelabs.org/classify.html}{detexify} 网站手写识别。另外还有数学公式识别软件 \href{https://mathpix.com/}{mathpix} 。

下面简单介绍一下定理、证明等环境的使用。
\begin{definition}
    定义环境
    \label{def:nosense}
\end{definition}
\cref{def:nosense}除了告诉你怎么使用这个环境以外，没有什么其它的意义。

除了 definition 环境，还可以使用 theorem 、lemma、corollary、assumption、conjecture、axiom、principle、problem、example、proof、solution 这些环境，根据论文的实际需求合理使用。

\begin{theorem}
    这是一个定理。
    \label{thm:example}
\end{theorem}
由\cref{thm:example}我们知道了定理环境的使用。

\begin{lemma}
    这是一个引理。
    \label{lem:example}
\end{lemma}
由\cref{lem:example}我们知道了引理环境的使用。

\begin{corollary}
    这是一个推论。
    \label{cor:example}
\end{corollary}
由\cref{cor:example}我们知道了推论环境的使用。

\begin{assumption}
    这是一个假设。
    \label{asu:example}
\end{assumption}
由\cref{asu:example}我们知道了假设环境的使用。

\begin{conjecture}
    这是一个猜想。
    \label{con:example}
\end{conjecture}
由\cref{con:example}我们知道了猜想环境的使用。

\begin{axiom}
    这是一个公理。
    \label{axi:example}
\end{axiom}
由\cref{axi:example}我们知道了公理环境的使用。

\begin{principle}
    这是一个定律。
    \label{pri:example}
\end{principle}
由\cref{pri:example}我们知道了定律环境的使用。

\begin{problem}
    这是一个问题。
    \label{pro:example}
\end{problem}
由\cref{pro:example}我们知道了问题环境的使用。

\begin{example}
    这是一个例子。
    \label{exa:example}
\end{example}
由\cref{exa:example}我们知道了例子环境的使用。

\begin{proof}
    这是一个证明。
    \label{prf:example}
\end{proof}
由\cref{prf:example}我们知道了证明环境的使用。

\begin{solution}
    这是一个解。
    \label{sol:example}
\end{solution}
由\cref{sol:example}我们知道了解环境的使用。
